% Proposed isolated A.90. No frozen29 source is changed.
\statement{A.90}{Proved}{The exact remaining odd-head coefficient criterion}
Assume the actual negative involution belongs to full
\(\widehat{GT}\), and retain A.89's notation
\(a=a_{\rm reg}\), \(d=d_\alpha\), \(\Theta_u\) and \(R_d\).
For every odd positive integer n let
\(\Lambda_n=\mathbb F_2[C_n]\), with projection denoted by
a subscript n. Then
\[
 (\Theta_uD_d)_n=1
       \quad\Longleftrightarrow\quad a\equiv d\pmod n.
 \tag{A.90.1}
\]
Consequently
\[
 \boxed{\ a_{\rm reg}=d_\alpha
 \ \Longleftrightarrow\
 (\Theta_uD_{d_\alpha})_n=1\text{ for every odd }n\ }
 \tag{A.90.2}
\]
and equivalently the remaining condition is
\[
             \overline{R_{d,n}}=T^dR_{d,n}
                    \quad\text{for every odd }n.
 \tag{A.90.3}
\]
These are exact criteria; no all-odd-head identity is asserted.
After A.89's dyadic equality, mixed even-times-odd heads impose
no additional exponent comparison.

\begin{proof}
A.88 gives \(\Theta_uD_a=1\), so \(\Theta_{u,n}\) is a unit.
The left side of (A.90.1) is therefore equivalent to
\(D_{d,n}=D_{a,n}\). Multiplication by \(T-1\) then gives
\(T^d=T^a\), which is equivalent to \(a=d\pmod n\) by the
group basis in \(\mathbb F_2[C_n]\).

Conversely suppose \(a=d\pmod n\). The difference
\(D_{d,n}+D_{a,n}\) is killed by \(T-1\). A vector fixed by
cyclic translation on the single regular orbit has constant
coefficients, so it has the form
\[
                    c(1+T+\cdots+T^{n-1}),\qquad c\in\mathbb F_2.
\]
Both a and d are full profinite units and hence odd. Their
geometric sums have augmentation one, so the displayed
difference has augmentation zero. The norm vector has
augmentation \(n\bmod2=1\). Hence \(c=0\), proving the
reverse implication. Here \(D_{b,n}\) is the projection of
the actual full continuous geometric sum; it is not defined
by division by \(T-1\) in the finite ring.

By A.89 the dyadic components of a and d already agree under
the stated full-GT hypothesis. Agreement modulo every odd n
is exactly agreement of all their odd-primary coordinates.
Together these are equality in \(\widehat{\mathbb Z}\).
Conversely that equality gives every displayed head identity.
This proves (A.90.2). A.89's exact hexagon-class identity gives
(A.90.3).

An equivalent character formulation, if desired, is
\[
 \Theta_u(\zeta)=\frac{\zeta-1}{\zeta^d-1}
       \quad\text{for every }\zeta\ne1
       \text{ of odd order in }\overline{\mathbb F}_2.
 \tag{A.90.4}
\]
For odd n, the polynomial \(X^n-1\) has distinct roots in
an algebraic closure of \(\mathbb F_2\), so evaluations at
these roots separate the finite cyclic ring. The denominator
in (A.90.4) is nonzero because d is a unit modulo the order
of \(\zeta\); the trivial character is already augmentation
one. This is only a restatement of the proved finite-ring
criterion, not an interpolation assumption.

A.42's prescribed pro-p hexagons for odd p do not directly
verify (A.90.3). These are relative mod-two kernel quotients
with odd cyclic retained heads, not maximal pro-p quotients.
An odd pro-p kernel has no nontrivial mod-two abelian
quotient. The dyadic module projection in A.89 therefore
cannot simply be repeated with an odd pro-p head.
For arbitrary \(\widehat{GT}_0\) data the dyadic premise of
(A.90.2) is not available; A.89's full-head criterion remains
the unconditional statement. In every case the full nonlinear
\(\Xi_d/\Omega_\alpha\) equation remains a separate gate.
\end{proof}
